\section{宋、辽与西夏之间：边界不是一条线}

从开封向北、西北看去，宋面对的不是两片静止的“少数民族边疆”，而是辽与西夏两个会征税、调兵、遣使、经营城市和书写自身合法性的国家。辽的多京、部族组织与农耕地区，西夏的河西走廊、灵州平原、城寨和寺院网络，各有资源与政治节奏。三方既交战也互市；同一条道路可以运兵、运盐、送使者，或成为逃亡者与商人的通道。把这段历史只写成宋朝的北方压力，便会把辽、西夏的主动选择从画面中抹去。

\subsection{两种边疆国家，两种资源组合}

耶律阿保机在十世纪初取得部族领导地位、继而建国，所形成的辽并非简单地由“游牧制度”加上一层汉式官名。契丹贵族、北方草原与东北地区的关系，多京城市、燕云农耕人口和南面行政的关系，须在不同地方分别考察。多中心的统治使辽能把草原机动、城市税源与南北交通连在一起；它也要求在皇族、后族、部族首领、地方官与宗教机构之间持续协调。所谓“南北面官”提供理解这一问题的入口，却不能替代对实际任官、税役和地域差异的考察。

党项李氏由定难军的地方势力发展为西夏王朝，走的又是另一条路。李元昊称帝并不是一支部落突然换上帝国名号，而是以河套、河西通道、灌溉绿洲、堡寨和宋辽之间的外交空间为条件的国家化。西夏文字、译经与法典工程显示朝廷将权威落到文书和宗教上的努力；黑水城出土的契约、账簿、佛经和行政文书，则让研究者看见官名之外的借贷、寺院、差役和地方书写。文书集中于某些地点、年代和保存条件，故不能替全境发言；但它确实纠正了仅靠宋人“夏国传”观看西夏的失衡。

\subsection{和平是一项有成本的制度}

澶渊盟约、宋夏庆历和议常被压缩为一方“纳币”或另一方“称臣”的胜负表。实际运行的是更复杂的制度：银绢的交付、使节的接待、边市的开放与限制、寨堡的修守、叛逃者的处置，以及各自国内如何解释对方的名号。对宋而言，支付可以减少一次全面动员的风险；对辽和西夏而言，它也是维持边境秩序与国内资源分配的一项收入和政治承认。没有哪一份盟书能取消军事侦察和局部冲突，正因如此，和平不能简单等同于没有战争。

它也不是没有代价的安定。宋廷须把银绢、军粮和河防长期纳入预算；辽须调和南部城市、防务与贵族联盟；西夏则在水源、通道与王室政治的压力下，处理与两大邻国的关系。永乐城等战事提醒我们，一座堡寨的得失会迅速改变粮道、声望与谈判位置，但不能从战报的夸张军数推出整个社会的动员规模。边疆的秩序是日常的行政劳动，不是一纸条约带来的静态地图。

\subsection{崩解、迁徙与材料的方向}

十二世纪女真兴起后，辽的既有联盟和军政网络受到冲击。辽亡不是某种制度“天然落后”的证明；不同地区、部族和城市在战争中作出不同选择，耶律大石西行亦使契丹政治网络延伸到中亚。西夏则在辽、金、蒙古连续重组的夹缝里调整外交，最终在蒙古扩张中覆亡。征服改变了城址、道路、人口与档案的保存，故“亡国”不能被理解为居民和制度在一日之间消失。

辽史的本纪和志书多由后来的元朝官修，宋人使辽记录带有外交观察者的视角；西夏的传世叙事又很大部分来自对手。与碑刻、考古和黑水城文书并读，并不会自动消除矛盾，却能把问题改得更具体：一条制度在哪些地方被实行？一个和约怎样被地方承受？哪些人没有进入朝廷的叙事？这也是本章把辽与西夏置于同一阅读路径的理由。
